% !PRINTED 662
\weberpage{662}
Machen wir diese Voraussetzung über die Functionale $\varphi_n$, so
können wir das \gesperrt{Kummer'sche} Theorem in der Fassung des
\S.\,174, 3. anwenden.

Wir bezeichnen in der Folge mit $\varepsilon$ Einheitsfunctionale, auf
deren genauere Kenntniss nichts ankommt, und setzen
\weq{7}{\Phi_n = \prod^{t} \varphi_{nt'}^{t} = \varepsilon\, \vartheta_n^m,}
wenn $t$ die durch $q$ nicht theilbaren Reste von $m$, die kleiner
als $m$ sind, durchläuft, und $t'$ durch $t\,t' \equiv 1\ (\mathrm{mod}\ m)$ bestimmt
ist. Dann ist $\vartheta_n$ eine Kreistheilungszahl in einem höheren
Körper, deren $m^{\text{te}}$ Potenz in $\Omega_m$ enthalten ist und die der zweiten
Bedingung genügt, dass
\weq{8}{\vartheta_n^{-1} \vartheta_1^n = \alpha}
eine Zahl des Körpers $\Omega_m$ ist.

Wir führen nun ein vielfach gebrauchtes Zeichen ein, indem
wir unter $E(x)$ die grösste ganze Zahl verstehen, die in der
reellen Zahl $x$ enthalten ist, und unter $(x)$ den Rest, der stets
ein echter Bruch ist, und, wenn $x$ selbst eine ganze Zahl sein
sollte, gleich Null zu setzen ist. Dann ist
\weq{9}{x = E(x) + (x).}

Ist $x$ eine ganze Zahl, so ist nach dieser Bezeichnungsweise
\[
m \left(\frac{x}{m}\right)
\]
auch eine ganze Zahl, und zwar der kleinste positive Rest der
Theilung von $x$ durch $m$.

In der Formel (7) ist der Index der Functionale $\varphi$ nur
nach dem Modul $m$ bestimmt. Der Exponent aber ist auf seinen
kleinsten Rest nach dem Modul $m$ reducirt. Ersetzen wir also
in (7) den Index $t'$ durch $n' t'$, so muss $t$ durch den Rest der
Division von $n t$ durch $m$, d.\,h. [nach (9)] durch
\[
m \left(\frac{n t}{m}\right) = n t - m E\left(\frac{n t}{m}\right)
\]
ersetzt werden, und wir können setzen
\[
\Phi_n = \prod^{t} \varphi_{t'}^{\,n t - m E\left(\frac{n t}{m}\right)} = \varepsilon\, \vartheta_n^m,
\]
andererseits ist
\[
\Phi_1 = \prod^{t} \varphi_{t'}^{t} = \varepsilon\, \vartheta_1^m,
\]
%CONT>>
